% ============================================
% DOC-07 (ZH)
% 中文翻译 — SPEC-935-R356
% ============================================
\section*{\textbf{DOC-07} --- 精神科诊断（1979）}
\markboth{DOC-07 --- 精神科诊断（1979）}{}

\noindent \textit{DSM-II —— 埃托尔·奥利维拉医生}

\noindent \textbf{患者:} 若阿金·安东尼奥·科雷亚 —— 登记号1,260
\textbf{医生:} 埃托尔·奥利维拉医生，CRM-MG 4,892
\textbf{日期:} 1979年4月14日

\vspace{0.5em}\noindent\rule{\textwidth}{0.5pt}\vspace{0.5em}

\noindent \textbf{CID-8 / DSM-II}

\noindent \textbf{轴I —— 主要诊断:}
- \textbf{295.3 —— 精神分裂症，偏执型}
  - 理由：系统化的被害妄想（相信一只“怪物”栖居其身）、持续存在的视觉与触觉幻觉、言语不连贯并带有新造词（“莫兰布多斯”作为一种存在状态）、情感迟钝与攻击性爆发交替出现。

\noindent \textbf{轴II —— 次要诊断:}
- \textbf{301.7 —— 病态人格}
  - 理由：无法感受悔恨、反复的攻击行为（青年时期三次确认的杀人）、操纵医疗人员、缺乏情感纽带。

\noindent \textbf{轴III —— 相关医学状况:}
- 慢性黄疸（右眼 —— 可能为童年严重营养不良后遗症，或未经治疗的肝病）。
- 蛋白质—热量营养不良与虚构时间线一致：1915年大旱，1917年转送。

\noindent \textbf{轴IV —— 心理社会应激源严重程度:}
- 极端。

\noindent \textbf{轴V —— 过去一年最高适应功能水平:}
- 10/100 —— 对自身与他人构成持续危险。

\vspace{0.5em}\noindent\rule{\textwidth}{0.5pt}\vspace{0.5em}

\noindent \textbf{临床观察:}

\noindent 患者于1917年、10岁时被收治入收容医院。历史记录显示，他抵达时“处于饥饿衰竭状态，来自塞阿拉的集中营，父母双亡，一只黄眼，除了‘饿’字之外一个字也不会说”。

\noindent 在62年的机构化收治期间，患者呈现：
- 反复自残发作（抓挠自己的脸直至流血）；
- 三次确认的杀人（1922、1937、1948）——所有受害者都是试图接近的工作人员或患者；
- 选择性缄默期（拒绝与工作人员说话；曾7年不与工作人员交谈——1949至1956年）；
- 强迫性书写一本“日记”，从未允许任何医生通读；
- 固着于“饥饿”和一只“黄眼怪物”的主题。

\noindent 患者表现出保存完好的智力——可能高于平均水平——但完全被导向其妄想的建构。他在妄想系统内部以内在连贯性叙述自己故事的能力令人瞩目。

\noindent \textbf{预后:} 谨慎。患者不可康复。维持隔离与持续镇静。

\vspace{0.5em}\noindent\rule{\textwidth}{0.5pt}\vspace{0.5em}

\noindent \textit{埃托尔·奥利维拉医生}
\textit{主治精神科医生 —— 最高安全病区}
\textit{巴尔巴塞纳收容医院}

\vspace{0.5em}
\noindent \textit{附言：埃托尔·奥利维拉医生于1979年6月13日失踪，就在签署本报告两个月之后。失踪案从未结案。在医院的人事档案里，一位同事在页边手写道：\textquotedblleft{}他在那之前三周就不吃东西了。他说自己不再感到饥饿——而他说，这是所有症状里最糟的一个。\textquotedblright{}}

\noindent\textit{\textbf{↪ 路线:} \rota{DOC-08} \rotasep \rota{DOC-04}}
